%% kapa: exam=1 round=27 year=2016 period=2 form=A subject=appraisal_law subject_ko=감정평가관계법규 seq=33 number=33
%% source: Q-Net 감정평가사 기출문제 (cst00309) artlSeq=5169402 file=extracted/1008/5169402/제27회 감정평가사 1차 문제/제27회 감정평가사 2교시문제(A형).pdf
%% conversion: kordoc 4.24.51 -> tools/gen1.py (auto); review=proofread-sonnet; flags=none
\begin{question}{33}
부동산등기법령상 등기절차에 관한 설명으로 옳지 않은 것은?
\begin{choices}
  \item 신탁재산에 속하는 부동산의 신탁등기는 수탁자가 단독으로 신청한다.
  \item 부동산표시의 변경등기는 소유권의 등기명의인이 단독으로 신청한다.
  \item 등기관이 등기의 착오가 등기관의 잘못으로 인한 것임을 발견한 경우 등기상 이해관계 있는 제3자가 없다면 지체 없이 그 등기를 직권으로 경정하여야 한다.
  \item 등기관은 등기권리자, 등기의무자 또는 등기명의인이 각 2인 이상인 경우에는 직권으로 경정등기를 한 사실을 그 모두에게 알려야 한다.
  \item 토지가 멸실된 경우에는 그 토지소유권의 등기명의인은 그 사실이 있는 때부터 1개월 이내에 그 등기를 신청하여야 한다.
\end{choices}
\end{question}
\answer{4}
